\chapter*{Abstract}
\addcontentsline{toc}{chapter}{Abstract}
\noindent

Sizing wind--photovoltaic--electrolyzer--battery systems is a black-box optimization problem: each design requires a sequential hourly simulation, and the constraint on the energy sold to the grid can only be verified after completing the annual horizon. This thesis develops a reproducible framework for comparing classic and hybrid metaheuristics on that problem, using the levelized cost of energy (LCOE) as the sole objective. Net present cost and the thesis-defined levelized cost of hydrogen are reported as secondary indicators.

First, the WPEB model by Li et al. \cite{Li2024WPEB} is reconstructed; absolute and relative errors support a partial reproduction rather than exact equivalence. The problem is then formulated as LCOE minimization over a mixed domain, with continuous wind capacity and discrete catalogs for the electrolyzer and the battery, and eight instances are defined by varying the grid-export limit, the catalog granularity and the plant scale. Each instance is also solved by brute force over the domain discretized at 0.01 MW, between \BFEDesignsEn{} and \BFCDesignsEn{} annual designs, which yields an exact optimum as the reference for every reported gap.

Seven algorithms are compared under an identical budget of \RequestsPerRun{} objective requests and \Replicates{} independent runs per instance: five classic metaheuristics (GWO, PSO, GA, SHO and WOA), one low-level integrative hybrid (PSO+GA), built without altering the implementation of its parent methods, and random search as the baseline. The comparison follows a seven-step statistical protocol with two-sided tests and a significance level and practical tolerance fixed a priori, which separates three questions: on a single instance, the Mann--Whitney test on the runs; in general, the Wilcoxon signed-rank test on one median per instance; and for all seven algorithms at once, the Friedman test with the Iman--Davenport correction. Differences are quantified with the $\hat{A}_{12}$ effect size and bootstrap intervals, and corrected with Holm's procedure.

The joint test rejects that the seven algorithms perform equally. By mean rank over the eight instances, \BestMeanRankAlgorithm{} and \SecondMeanRankAlgorithm{} take the first two positions and the hybrid ranks \ControlRankOrdinalEn{} out of seven, between its two parents. Instance by instance, PSO+GA is not distinguishable from PSO on any of the eight instances, is significantly better than GA on \CtrlGaPlus{} of them and outperforms random search on all eight; overall it trails PSO by a small margin (adjusted $p$-value of \CtrlPsoWilcoxonHolmEn) and is not distinguishable from GA (\CtrlGaWilcoxonHolmEn). With this configuration and budget, hybridization does not improve on the better of its parents. The convergence analysis qualifies this result: the hybrid has the best mean rank up to \ConvControlLeadTo{} requests and stagnates afterwards, and the final order only settles from \ConvSettleFrom{} requests onwards. The executed notebook generates the protocol outputs, tables and figures, making budgets, seeds, duplicates, feasibility and every reported result auditable.

\vspace{1em}
\noindent\textbf{Keywords:} hybrid metaheuristics; statistical comparison of algorithms; green hydrogen; sizing; levelized cost of energy; empirical algorithm evaluation.
